\section{Full sources, endpoint observables, and certified heat: Rounds 23--26}
\label{sec:recent}

The twenty-eight investigations S1--AF2 develop three separate models.  The
homogeneous initial-diagonal problem uses dimensionless generators with physical
energy unit $\alpha/8$; the finite-graph problem uses $L=H/\alpha$ and heat time
$\sigma=\alpha t/\hbar$; the canonical summable family uses its original real-time
generator and a $q$-dependent interaction profile.  Throughout, lattice spacing
$a$, $E_\star>0$, $\alpha/E_\star>0$, and $\hbar>0$ are fixed.  A source filter,
collar depth, Galerkin projection, joining time, or input-overlap restriction is a
proof or approximation parameter, not a new term in an action.

The marker \workstar\ means a contribution newly derived or computed in this
workbench; scientific priority is unverified.  S1/S2/T1, all ten Round 24 loops,
and all ten Round 26 loops have separately authored model-agent derivations and
skeptical reviews.  T2/U1/U2 and AA1/AA2 instead have same-author derivation,
reverse reconstruction, and critical self-review.  These categories describe
provenance, not external mathematical peer review or experimental independence.
Historical limited verdicts remain limited even where later work answers a
previously missing premise.  The detailed source-to-equation mapping is recorded
in \repo{audit/recent-ledger.json} accompanying this manuscript.

\subsection{An actual homogeneous source and its regulated inverses}
\label{recent:homogeneous}

\subsubsection{\workstar\ S1: a generated cubic source and the complete crossing defect}

Use the complete 24-link onsite factors and four-site stars
$Y=\{0,e_1,e_2,e_3\}$ of the homogeneous construction.  On one star let $H=H_{0,Y}$,
$P_0=|\Omega_Y\rangle\langle\Omega_Y|$, $Q_0=I-P_0$, and $H\ge Q_0$.  With the
actual 21-face potential $\phi_0$, set
\begin{equation}
 v=\phi_0\Omega_Y,\quad u=H_{Q_0}^{-1}v,\quad
 a_0=\langle u,v\rangle,\quad b_0=\norm{u}^2,\quad
 s_0^2=\norm{v}^2=\frac7{12}\tau^2 .
 \label{recent:source-data}
\end{equation}
The Haar identity follows because each face has a free Haar link and distinct
faces have an unmatched free link: the 21 diagonal half-trace moments are
$1/4$ and every cross term vanishes.  Define
$S=|u\rangle\langle\Omega_Y|-|\Omega_Y\rangle\langle u|$ and
$A_0=|v\rangle\langle\Omega_Y|+\mathrm{adjoint}$.  The specified repeated-anchor
term in the previous remainder is
\begin{equation}
 C=\frac12[S,[S,\phi_0]]-\frac16[S,[S,A_0]],\qquad
 w=Q_0C\Omega_Y=-a_0u-\frac{b_0}{3}v. \qquad\text{\workstar}
 \label{recent:cubic-source}
\end{equation}
The two double commutators have the same $Q_0$-column
$-3a_0u-b_0v$; their coefficients give the displayed result.  In particular,
$\langle u,w\rangle=-4a_0b_0/3<0$ for $\tau\ne0$, so this is an actual nonzero
source, with $r:=\norm{w}\le4s_0^3/3$.  It is one indexed cubic term; other
terms in a regrouped cubic coefficient could cancel it.  Its exact local split
is $C=c_0I+A+Q_0(C-c_0I)Q_0$, where $c_0=\langle\Omega_Y,C\Omega_Y\rangle$
and $A=|w\rangle\langle\Omega_Y|+\mathrm{adjoint}$.
All relative ratios divided by $r$ below require $\tau\ne0$, hence $r>0$.
At $\tau=0$ the generated source vanishes; the undivided bounds still apply.

Write $M=7|\tau|\le35/1664$, $D_c=Q_c\phi_cQ_c$, and
$G_\Lambda=H_{0,\Lambda}+\sum_cD_c$, retaining only complete stars in a containing
finite cuboid.  The one-star operator $G_Y=H+D_0$ satisfies
$(1-M)H\le G_Y\le(1+M)H$ as forms and
$D(G_Y)=D(H)$ by bounded perturbation.  Therefore
\begin{equation}
 z=(G_Y|_{Q_0})^{-1}w,\quad
 K_Y=|z\rangle\langle\Omega_Y|-\mathrm{adjoint},\quad
 [K_Y,G_Y]=-A,\qquad \norm{K_Y}\le\frac{r}{1-M}.
 \label{recent:interior-inverse}
\end{equation}
The inverse places $z$ in the operator domain.  Identity extension preserves
$D(H_{0,\Lambda})$ because the local commutator with $H$ is bounded and exterior
spectral projections commute with $K_Y$.  The graph-bounded exponential and its
inverse preserve that domain as well.  The complete extended equation is
\begin{equation}
 [K_Y\otimes I,G_\Lambda]=-A\otimes I+F_Y,\qquad
 F_Y=\sum_{c\text{ crossing }Y}[K_Y,D_c],\qquad
 \norm{F_Y}\le\frac{6Mr}{1-M}. \qquad\text{\workstar}
 \label{recent:crossing-defect}
\end{equation}
At the positive-octant origin there are three possible crossings, with anchors
$e_1,e_2,e_3$, each on a seven-site union.  For the single-source decomposition,
$\sum_c e^{7\mu}\norm{[K_Y,D_c]}\le6e^{7\mu}Mr/(1-M)$.  A translated interior
star has twelve crossings.  The reverse-origin sharper rooted count is
$12\cdot7=84$, giving a family bound $168e^{7\mu}Mr_\star/(1-M)$; the forward
count 96 gives the valid weaker constant 192.  These are one-step interaction
bounds, not an all-order inverse.

An exterior mean-zero unit Haar character $\xi$ distinguishes the actual source
from its vacuum-projected replacement:
$(A\otimes I)(\Omega_Y\otimes\xi)=w\otimes\xi$ while
$(A\otimes P_{\mathrm{ext}})(\Omega_Y\otimes\xi)=0$.  Also a bounded
$D(G)$-preserving exact inverse must satisfy $P_EAP_E=0$ at every energy; a
positive ground gap alone says nothing about all excited Bohr differences.
This explains S1's limited gate and motivates a regulated full-source equation.
Sources: \repo{research/round23/forward/s1/report.md},
\repo{research/round23/reverse/s1/report.md}, and
\repo{research/round23/skeptic/s1.md}.

\subsubsection{\workstar\ S2 and W1--W2: retain the residual and name its topology}

For $\alpha_s^G(A)=e^{isG}Ae^{-isG}$, S2 fixes $0<\Theta\le1/16$ and defines
\begin{align}
 \mathcal L_\Theta(A)&=-\frac i2\int_{-\Theta}^{\Theta}
 \operatorname{sgn}(s)(1-|s|/\Theta)\alpha_s^G(A)\,ds,\notag\\
 \mathcal R_\Theta(A)&=\frac1{2\Theta}\int_{-\Theta}^{\Theta}
 \alpha_s^G(A)\,ds,\qquad
 [\mathcal L_\Theta(A),G]=-A+\mathcal R_\Theta(A). \qquad\text{\workstar}
 \label{recent:short-filter}
\end{align}
These are vectorwise strong integrals.  In weak pairings on $D(G)$, integration
by parts uses the kernel derivative
$2\delta_0-\Theta^{-1}\mathbf1_{[-\Theta,\Theta]}$.  The bounded weak
commutator then proves $\mathcal L_\Theta(A)D(G)\subset D(G)$ by self-adjointness;
domain preservation of an arbitrary initial $A$ is not assumed.  Its norm and
graph bounds are $\Theta r/2$ and $\Theta r/2+2r$.  At Bohr frequency $\omega$,
\begin{equation}
 \ell_\Theta(\omega)=\frac{1-\operatorname{sinc}(\Theta\omega)}{\omega},
 \qquad r_\Theta(\omega)=\operatorname{sinc}(\Theta\omega),
 \qquad \ell_\Theta(0)=0, r_\Theta(0)=1.
 \label{recent:sinc-multipliers}
\end{equation}
Thus every equal-energy block survives exactly; shortening $\Theta$ leaves more
of the source.

For the rooted norm
$\norm{\Phi}_2=\sup_x\sum_{X\ni x}2^{|X|}\norm{\Phi_X}$, a surviving ordered
commutator word starts on one star and adds at most three sites at each step.
The forward rooted count, before simplex integration, is
\begin{equation}
 N_0\le64r_\star,\quad N_{n+1}\le64M(8+3n)N_n,
 \quad N_n\le64r_\star(192M)^n(8/3)_n. \qquad\text{\workstar}
 \label{recent:root-count}
\end{equation}
Here the insertion cost includes $2M$ from a commutator and $2^3$ from the
support weight; incoming anchors and repetitions remain.  The simplex contributes
$|s|^n/n!$, and the inverse and residual filter moments are respectively
$\Theta^{n+1}/(n+2)!$ and $\Theta^n/(n+1)!$.  Consequently
\begin{equation}
 \norm{\mathcal L_\Theta(\{A_b\})}_2
 \le32\Theta r_\star(1-192M\Theta)^{-8/3},\qquad
 \norm{\mathcal R_\Theta(\{A_b\})}_2
 \le64r_\star(1-192M\Theta)^{-8/3},
 \label{recent:s2-weighted}
\end{equation}
with $192M\Theta\le105/416<1$.  Independently, reverse counts at most
$13^nn!$ relative words and $4+3n$ rooted translates.  With $u=208M\Theta$,
its orbital sum is $16r_\star(4-u)/(1-u)^2$; the inverse and residual bounds
are $8\Theta r_\star(4-u)/(1-u)^2$ and $16r_\star(4-u)/(1-u)^2$.
At $u=35/128$ the rational sum is $6784/961$.  The full tail after depth $N$ is
\begin{equation}
 \sum_{n>N}(4+3n)u^n
 =\frac{u^{N+1}[(3N+7)-(3N+4)u]}{(1-u)^2}.
 \label{recent:s2-tail}
\end{equation}
These alternative bounds concern indexed interaction decompositions, not a
bounded extensive infinite sum.

The forward-origin infinite-representation identification adds a separate
form argument.  In the admitted product-reference representation, write
$G\leftrightarrow h_0+d$, $\kappa=4M<1$, and
$\sum_b\norm{Q_bu}^2\le4h_0[u]$.  Subtracting resolvent form equations for
$u=(G+1)^{-1}f$ and $u_\Lambda=(G_\Lambda+1)^{-1}f$ gives
\begin{equation}
 \norm{u_\Lambda-u}_{\mathcal F}
 \le\frac{\sqrt\kappa}{1-\kappa}
 \left(M\sum_{b\text{ omitted}}\norm{Q_bu}^2\right)^{1/2}\longrightarrow0.
 \label{recent:form-identification}
\end{equation}
Strong-resolvent convergence identifies the strong unitary limit.  Stabilization
of every finite connected word plus the uniform summable tail improves the
single-source short-time limit to operator norm.  It does not prove
$D(G)=D(H_0)$ globally or a norm limit of global propagators.
Sources: \repo{research/round23/forward/s2/report.md},
\repo{research/round23/reverse/s2/report.md},
\repo{research/round23/skeptic/s2.md}.

W1 then evaluates the actual source's first form moment, a reverse-origin
refinement separately reviewed.  The ground-transform identity and the
half-trace gradient identity give
\begin{equation}
 h[F\Omega_Y]=8\sum_e\int |\nabla_eF|^2|\Omega_Y|^2,
 \quad\sum_a|X_e^aW_f|^2=\frac{1-W_f^2}{4},\quad
 h[v]=14\tau^2=24s_0^2. \qquad\text{\workstar}
 \label{recent:source-moment}
\end{equation}
Free-link parity also cancels cross gradients.  In the spectral measure of $v$,
$w=-f(H)v$ with decreasing positive $f(\lambda)=a_0/\lambda+b_0/3$.
The covariance of $\lambda$ with $f(\lambda)^2$ is nonpositive, so
$h[w]/r^2\le24$; Jensen's inequalities for the two inverse moments yield
$r\ge s_0^3/432$.  Because $G\Omega=0$,
$\mathcal R_\Theta(A)\Omega=\operatorname{sinc}(\Theta G)w$.
Applying $\operatorname{sinc}x\ge1-x/2$ to this positive spectral measure and
$G\le(1+4M)H_0$ gives the actual lower bound
\begin{equation}
 \norm{\mathcal R_\Theta(A)}
 \ge[1-12\Theta(1+4M)]r\ge\frac{311}{1664}r>0.
 \label{recent:residual-lower}
\end{equation}
This prevents deleting the short-filter cubic residual.  It does not establish
an actual zero-frequency obstruction to every inverse.  The complementary
forward result is the source-mass certificate
$\norm{\mathcal R_\Theta(A)}\ge(5/6)r\sqrt{\nu([0,E])}$ when $\Theta E\le1$,
and $\norm{\mathcal R_\Theta(A)}/r\to1$ as $\Theta\downarrow0$; neither is a
full-operator upper estimate.  If a bounded exact inverse $K$ were already
known, integration would instead give
$\norm{\mathcal R_\Theta(A)}\le\norm K/\Theta$, a conditional criterion only.
Sources: \repo{research/round24/reverse/w1/report.md},
\repo{research/round24/forward/w1/report.md},
\repo{research/round24/skeptic/w1-review.md}.

W2 holds $G$ fixed and sets $A_n=\mathcal R_\Theta^n(A)$,
$K_n=\sum_{j<n}\mathcal L_\Theta(A_j)$.  Telescoping preserves
\begin{equation}
 [K_n,G]=-A+A_n,\qquad
 \norm{K_n}\le n\Theta r/2,\qquad
 (A_n)_\omega=\operatorname{sinc}(\Theta\omega)^nA_\omega.
 \label{recent:iterate-filter}
\end{equation}
On the actual vacuum column, $g=1-4M\ge381/416$ supplies a gap.  With the
reverse bound $\rho=1-(\Theta g)^2/8<1$, a piecewise Taylor and
$|\operatorname{sinc}x|\le1/|x|$ argument yields
\begin{align}
 \norm{A_n\Omega}&\le\rho^nr,\notag\\
 \norm{K_n\Omega-G_Q^{-1}w}&\le\rho^nr/g,\qquad
 \norm{G(K_n\Omega-G_Q^{-1}w)}\le\rho^nr. \qquad\text{\workstar}
 \label{recent:column-convergence}
\end{align}
Forward independently obtains denominator 10 in $1-(\Theta g)^2/10$.
These are graph-norm convergence of a column, uniform in the containing volumes
and valid in the admitted infinite representation.  In each fixed finite volume,
compact resolvent and dominated convergence in an energy eigenbasis additionally
prove
\begin{equation}
 \mathcal R_\Theta^n(A)\xrightarrow{\mathrm{strong}}
 \mathcal D_G(A):=\sum_EP_EAP_E.
 \label{recent:pinching-limit}
\end{equation}
The value is not known to vanish.  Compact-resolvent doublets with splittings
$1/(k+1)$ and unit swap sources show why strong convergence can coexist with
residual norm one and an unbounded formal inverse.  In the actual model the old
weighted majorants cover only $n\le3$ at the cap; failure at $n=4$ is failure of
that estimate, not divergence of the dynamics.  W2 therefore does not close a
changed-diagonal nonlinear iteration.
Sources: \repo{research/round24/forward/w2/report.md},
\repo{research/round24/reverse/w2/report.md},
\repo{research/round24/skeptic/w2-review.md}.

\subsubsection{\workstar\ AB1--AB2: the surviving block is controlled by its boundary}

Pinching the complete identity \eqref{recent:crossing-defect} gives
\begin{equation}
 \mathcal D_G(A)=\mathcal D_G(F_Y),\qquad
 \norm{\mathcal D_G(A)}\le\frac{6Mr}{1-M},\qquad
 \frac{6M}{1-M}\le\frac{70}{543}. \qquad\text{\workstar}
 \label{recent:ab-resonance}
\end{equation}
The diagonal map is contractive because its input and output blocks are
orthogonal.  As $r=O(|\tau|^3)$, this is an actual uniform finite-volume quartic
upper bound, not an evaluation of a nonzero resonant block.  On a spectral window
$P=\mathbf1_{[E-\varepsilon,E+\varepsilon]}(G)$, reverse additionally obtains
\begin{equation}
 \norm{PAP}\le\frac{(2\varepsilon+6M)r}{1-M},
 \label{recent:ab-window}
\end{equation}
because $\norm{(G-E)P}\le\varepsilon$.  If no crossing star is retained,
$F_Y=0$ and the interior inverse is exact; otherwise zero versus nonzero crossed
blocks remains unresolved.  Pinching need not preserve locality.

For the Gaussian density $p_s(t)=e^{-t^2/(2s^2)}/(\sqrt{2\pi}s)$ and odd tail
$h_s(t)=\operatorname{sgn}(t)\int_{|t|}^\infty p_s(u)\,du$, define
$\mathcal R_s(A)=\int p_s(t)\alpha_t(A)dt$ and
$\mathcal J_s(A)=-i\int h_s(t)\alpha_t(A)dt$.  Since
$h_s'=\delta_0-p_s$, $\norm{h_s}_1=s\sqrt{2/\pi}$, and
$\norm{p_s'}_1=\sqrt{2/\pi}/s$, weak integration by parts proves
\begin{align}
 [\mathcal J_s(A),G]&=-A+\mathcal R_s(A),\qquad
 \mathcal J_s(A)D(G)\subset D(G),\notag\\
 \mathcal R_s(A)&=\mathcal R_s(F_Y)+i\int p_s'(t)\alpha_t(K_Y)\,dt,\notag\\
 \frac{\norm{\mathcal R_s(A)}}r
 &\le\min\left\{1,\frac{6M+\sqrt{2/\pi}/s}{1-M}\right\}. \qquad\text{\workstar}
 \label{recent:gaussian-contraction}
\end{align}
The residual and inverse spectral multipliers are
$e^{-s^2\omega^2/2}$ and $(1-e^{-s^2\omega^2/2})/\omega$, retaining the
zero-frequency block.  At $s=4$ the bound is $626/1629<0.385$; reverse's $s=2$
gives $1042/1629<2/3$.  This advances W2 from a column contraction to a
full-operator contraction on the fixed source and initial $G$.  The asymptotic
boundary budget is an upper bound, not a positive residual lower bound.
Gaussian time support is unbounded, so the old finite-radius connected series
cannot establish its locality by direct integration.
Sources: \repo{research/round26/forward/ab1/report.md},
\repo{research/round26/reverse/ab1/report.md},
\repo{research/round26/forward/ab2/report.md},
\repo{research/round26/skeptic/ab1.md},
\repo{research/round26/skeptic/ab2.md}.

\subsubsection{\workstar\ AE1--AE2: locality, graph error, and the original weighted norm}

AE1 uses an all-time interaction-picture commutator recurrence.  Each site lies
in at most four stars and a four-site star meets at most sixteen possible stars.
Removing internal unitaries in the commutator recurrence leaves chains, rather
than the growing-union count in \eqref{recent:root-count}.  For the forward
collar $C_R=\{x\in\mathbb N_0^3:|x|_1\le2R+1\}$, every chain of at most $R$
steps is retained.  Duhamel and the chain exponential then give
\begin{align}
 \norm{\alpha_t^G(A)-\alpha_t^{G_R}(A)}
 &\le r\sum_{m\ge R+1}\frac{(32M|t|)^m}{m!}
 \le r2^{-R-1}e^{64M|t|},\notag\\
 \norm{\mathcal R_s^G(A)-\mathcal R_s^{G_R}(A)}
 &\le r2^{-R}e^{(64Ms)^2/2}. \qquad\text{\workstar}
 \label{recent:gaussian-locality}
\end{align}
The last step integrates the entire Gaussian using
$\mathbb E e^{v|T|}\le2e^{v^2s^2/2}$.  At the old cap and $s=4$, $R=64$ gives
relative operator error below $2^{-42}$ but requires 374,660 complete factors.
No local finite Hilbert dimension is assumed.

The independent reverse construction uses $C_R=\{0,\ldots,R+1\}^3$ and
$b_R=3R^2+9R+7$ crossing stars.  With $v=64M$, it bounds the residual and inverse
errors by
\begin{align}
 \varepsilon_R&=4rMb_Rs(1+vs)e^{v^2s^2/2}2^{-R},\notag\\
 \kappa_R&=4rMb_Rs^2e^{v^2s^2/2}2^{-R},\qquad
 \norm{[\mathcal J_s-\mathcal J_{s,R},G]}\le2\varepsilon_R.
 \label{recent:gaussian-graph}
\end{align}
The graph conclusion is an additional calculation of every crossing commutator;
small operator-norm error by itself would not prove it.  At $s=2,R=40$, the
bounds are $\varepsilon_R/r<2\cdot10^{-7}$ and $\kappa_R/r<10^{-7}$ on
74,088 factors.  Telescoping collars produces exponential diameter decay, with
polynomial translation multiplicities harmless.  It does not establish the
old support-cardinality norm: a factor $e^{\mu O(R^3)}$ overwhelms this
$2^{-R}$ upper majorant.  That failed majorant is not a divergence theorem.
Sources: \repo{research/round26/forward/ae1/report.md},
\repo{research/round26/reverse/ae1/report.md},
\repo{research/round26/skeptic/ae1.md}.

AE2 obtains a simultaneous statement in the original rooted weight by explicitly
restricting $M\le10^{-3}$, equivalently $|\tau|\le1/7000$.  At the fixed proof
duration $T=3$, use
\begin{equation}
 p_T(t)=\frac{(1-|t|/T)_+}{T},\quad
 h_T(t)=\frac{\operatorname{sgn}(t)}2(1-|t|/T)_+^2,
 \quad \widehat p_T(\omega)=\operatorname{sinc}^2(T\omega/2).
 \label{recent:compact-kernel}
\end{equation}
Its derivative and moments are $h_T'=\delta_0-p_T$,
$\norm{p_T'}_1=2/T$, $\norm{h_T}_1=T/3$.  For a translated interior source,
all twelve crossings cost $24Mr_b/(1-M)$, so the boundary identity proves
\begin{equation}
 \norm{\mathcal R_T(A_b)}
 \le\frac{24M+2/T}{1-M}r_b
 \le\frac{2072}{2997}r_b<0.7r_b. \qquad\text{\workstar}
 \label{recent:compact-contraction}
\end{equation}
Origin-only three-crossing constants cannot be used for this family.
Integrating the same ordered words as S2 yields, with $c=192M$, $p=8/3$,
\begin{align}
 \norm{\mathcal R_T(\{A_b\})}_2
 &\le128r_\star\sum_{n\ge0}\frac{(p)_n(cT)^n}{(n+2)!},\notag\\
 \norm{\mathcal J_T(\{A_b\})}_2
 &\le128Tr_\star\sum_{n\ge0}\frac{(p)_n(cT)^n}{(n+3)!}.
 \label{recent:compact-weighted}
\end{align}
Here $cT\le72/125<1$, giving the common bound $64(125/53)^3r_\star$.
Reverse's independent $208M$ count gives the smaller displayed sufficient
budget $844000r_\star/2209<383r_\star$ at $208MT=78/125$.
These prove finiteness, not weighted contraction: their positive residual
series already begins at the initial $64r_\star$ budget.

The same certificates require simultaneously
$T>2/(1-25M)$ and $T<1/(192M)$, possible only if $409M<1$; the reverse radius
requires $441M<1$.  At the old cap $M=35/1664$, the first lower bound is
$3328/789$, while the upper bounds are $26/105$ and $8/35$.  This is a precise
failure of those joint certificates at the old cap.  At their radius endpoint
the inverse majorant is integrable but the residual majorant is not.  The next
missing premise is weighted contraction for generated sources with all scalar,
diagonal, and domain updates, not another application of the original-source
factor 0.7.
Sources: \repo{research/round26/forward/ae2/report.md},
\repo{research/round26/reverse/ae2/report.md},
\repo{research/round26/skeptic/ae2.md}.

\subsection{A fixed physical graph: from memory to computed relative heat}
\label{recent:finite-graph}

The graph throughout this subsection has vertices
$\{0,1,2\}\times\{0,1,2\}\times\{0,1\}$, 33 links, 20 square faces, and all
18 vertex Gauss constraints.  Its full physical Hilbert space is
$\mathcal H_{\rm phys}=L^2(\mathrm{SU}(2)^{33},d\mathrm{Haar})^{\mathrm{SU}(2)^{18}}$.
Define
\begin{equation}
 \begin{aligned}
 K&=\sum_e C_e,\qquad L=K+\lambda(20-S),\qquad
 S=\sum_{p=1}^{20}W_p,\\
 W_p&=\tfrac12\Tr U_{\partial p},\qquad
 0\le\lambda\le10^{-2},\qquad\sigma=\alpha t_{\rm phys}/\hbar.
 \end{aligned}
 \label{recent:graph-model}
\end{equation}
In heat estimates below, integration variables $t,s$ and joining times are dimensionless coordinates in the same units as $\sigma$; physical duration is $t_{\rm phys}=(\hbar/\alpha)\sigma$.
The multiplication $20-S$ lies between zero and 40.  Gauge averaging commutes
with the electric spectral projections; bounded perturbation preserves the
self-adjoint electric domain and its smooth invariant core.  The electric
resolvent is compact.  Its vacuum is unique and the physical electric gap is
3: a nonempty invariant active-edge support has no degree-one vertex, hence
contains a cycle, whose at least four spin-half edges cost $4(3/4)$.
A fundamental square character attains 3.  These are full-space facts, not a
finite-spin assumption.

\subsubsection{\workstar\ T1--T2: a full-space memory model and separate ground centering}

For the inherited infinite-rank isometry $J$ from the selected physical space,
write the full Hamiltonian in retained/complementary blocks
\begin{equation}
 H_\lambda=\begin{pmatrix}A&B^*\\B&C\end{pmatrix},\qquad
 \widetilde H_\lambda=\begin{pmatrix}A&B^*\\B&C_0\end{pmatrix},\quad
 C-C_0=\alpha\lambda D,\quad 0\le D\le40Q,
 \quad\norm B=\tfrac52\alpha\lambda.
 \label{recent:t-blocks}
\end{equation}
The diagonals have their actual inherited electric domains,
$A\ge19\alpha\lambda$ and $C_0\ge3\alpha$.
The scalar Schur/Young inequality gives
$\widetilde H_\lambda\ge18\alpha\lambda I$ throughout the stated interval;
$H_\lambda\ge\widetilde H_\lambda$ as forms gives the same separate spectral
floor for $H_\lambda$.  No order between their heat operators is inferred.
Bounded perturbation establishes both on the full electric domain.

Eliminating the complementary component shows that
$\widetilde T(t)=J^*e^{-t\widetilde H_\lambda/\hbar}J$ satisfies
\begin{align}
 \widetilde T(t)&=e^{-tA/\hbar}+\hbar^{-2}
 \int_{0\le u\le s\le t}e^{-(t-s)A/\hbar}
 K_2(s-u)\widetilde T(u)\,du\,ds,\notag\\
 K_2(v)&=B^*e^{-vC_0/\hbar}B. \qquad\text{\workstar}
 \label{recent:memory-return}
\end{align}
The full returned $\widetilde T(u)$ retains repeated excursions.  Iterating a
difference of two solutions bounds it by
$M_T(\norm Bt/\hbar)^{2n}/(2n)!$, proving uniqueness at each finite window.
The corresponding Schur complement has minus sign
$A+z-B^*(C_0+z)^{-1}B$.  A one-return substitution loses a fourth-order
$(B^*B)^2$ term.  The two earlier diagnostic channels cannot be made autonomous:
their full second-derivative defect dominates
$\alpha^2\lambda^2\left(\begin{smallmatrix}19/4&1/4\\1/4&19/4\end{smallmatrix}\right)>0$.
This last actual-model clarification originated in the independent T1 review.

Let $d=3-18\lambda$.  Duhamel for the purely complementary perturbation has
one leakage factor on each side, each bounded by
$(5\lambda/2)e^{-18\lambda\sigma}(1-e^{-d\sigma})/d$.  Thus
\begin{align}
 \norm{T(t)-\widetilde T(t)}
 &\le250\lambda^3e^{-18\lambda\sigma}F_d(\sigma),\notag\\
 F_d(\sigma)&=\frac{\sigma(1+e^{-d\sigma})-2(1-e^{-d\sigma})/d}{d^2}
 \le\min\{\sigma^3/6,\sigma/d^2\},\notag\\
 \sup_{t\ge0}\norm{T(t)-\widetilde T(t)}
 &\le\frac{125\lambda^2}{9e(3-18\lambda)^2}
 \le\frac{312500}{477144}\lambda^2.\qquad\text{\workstar}
 \label{recent:t-unshifted}
\end{align}
This is cubic in coupling at fixed $\sigma$, but its uniform quadratic bound
uses decay in the unshifted energy convention.

T2 answers the centering objection rather than multiplying this estimate by
uncontrolled ground exponentials.  Along
$L_s=\widetilde H_\lambda/\alpha+s\lambda D$,
min-max gives a simple ground $\epsilon_s\le40\lambda$ and excitation gap
$g=3-85\lambda/2\ge103/40$.  The complementary ground inverse has denominator
$d_0=3-40\lambda\ge13/5$, so its leakage is
$\ell=(5\lambda/2)/d_0$.  Differentiation gives
$\epsilon_s'=\lambda\langle\psi_s,D\psi_s\rangle$ and
$\norm{G_s'}\le40\lambda\ell/g$.  In centered Duhamel the ground--ground term
cancels exactly because
$G_s(\lambda D-\epsilon_s'I)G_s=0$.  Two mixed, magnetic excited, and scalar
excited terms cost respectively $80\lambda\ell/g$, $80\lambda\ell/g$, and
$40\lambda\ell^2/g$.  Integration in $s$ proves
\begin{align}
 \sup_{t\ge0}\norm{J^*e^{-t(H_\lambda-e_\lambda)/\hbar}J
 -J^*e^{-t(\widetilde H_\lambda-\widetilde e_\lambda)/\hbar}J}
 &\le60\lambda^2,\notag\\
 0\le e_\lambda-\widetilde e_\lambda
 &\le\tfrac{6250}{169}\alpha\lambda^3,\qquad
 \norm{G_\lambda-\widetilde G_\lambda}
 \le\tfrac{20000}{1339}\lambda^2.\qquad\text{\workstar}
 \label{recent:t-centered}
\end{align}
The exact cap coefficient before rounding is $1042500/17407<60$.
The excited remainder has initial complementary component $-QG_sJ$, which must
remain in the bound $\norm{Q(e^{-u(L_s-\epsilon_s)}-G_s)J}
\le\min\{2\ell,e^{-gu}\}$.  These results concern infinite-rank selected heat,
not a finite algorithm or a relative-output theorem.  T2 is same-author
self-review.  Sources: \repo{research/round23/forward/t1/report.md},
\repo{research/round23/reverse/t1/report.md},
\repo{research/round23/skeptic/t1.md},
\repo{research/round23/forward/t2/report.md},
\repo{research/round23/skeptic/t2.md}.

\subsubsection{\workstar\ X1: two truncation theorems and an actually assembled sector}

A zero-extended finite-rank heat operator has full-space error exactly one at
$\sigma=0$.  X1 therefore separates two meaningful questions.  For the kinetic
cutoff $P_R=\mathbf1_{[0,R)}(K)$, $R>3$, let $A_R=P_RLP_R$,
$B_R=Q_RLP_R$, $b=20\lambda$, and $c=20\lambda$.
Then $\norm{B_R}\le b$ and $Q_RLQ_R\ge R$.  Complementary variation of constants
and its return give
\begin{equation}
 \norm{e^{-\sigma L}-e^{-\sigma A_R}P_R}
 \le e^{-R\sigma}+\frac{2b}{R}+\frac{b^2\sigma}{R}.
 \label{recent:x-cutoff}
\end{equation}
For actual ground energies $\epsilon,\mu_R$, Schur elimination and the full
excited threshold 3 imply
\begin{equation}
 0\le\mu_R-\epsilon\le\frac{b^2}{R-c},\qquad
 \norm{G-G_R}\le D_R:=\frac{b}{R-c}\left(1+\frac b{3-c}\right).
 \label{recent:x-cutoff-ground}
\end{equation}
For any $0<\sigma_0\le T$, separate ground centering consequently gives
\begin{align}
 \sup_{\sigma\ge\sigma_0}
 \norm{e^{-\sigma(L-\epsilon)}-e^{-\sigma(A_R-\mu_R)}P_R}
 &\le\min\{2,\max(F_R,D_R+2e^{-(3-c)T})\},\notag\\
 F_R&=e^{cT}\left[e^{-R\sigma_0}+2b/R+b^2T/R+Tb^2/(R-c)\right].
 \qquad\text{\workstar}
 \label{recent:x-delayed}
\end{align}
Choosing $T$ first and then $R$ proves uniform operator-norm convergence on every
positive delayed half-line.  The certified example $\lambda\le0.01$,
$\sigma_0=1,T=4,R=1024$ has error below 0.002; that large matrix was not assembled.
A Hermitian finite coefficient perturbation of norm $\zeta$ adds at most
$4\zeta/(g-2\zeta)$ uniformly when both finite matrices use their own grounds
and $2\zeta<g$.  Independent rounding of a ground exponent instead produces
$e^{\sigma\delta}$ growth or loss of the limiting ground.

The independent reverse projection retains all invariant Peter--Weyl blocks
with degree $d=\sum_e2j_e\le N$.  A face multiplier raises degree by at most four.
For initial degree $d_0\le N$, set $r=\lfloor(N-d_0)/4\rfloor+1$ and
$z=20\lambda\sigma$.  Dyson words of order below $r$ agree exactly, including
all intermediate states; therefore their uncentered difference is at most
$2e^{-z}\sum_{n\ge r}z^n/n!\le2z^r/r!$.
The true ground Neumann series has ratio $a=20\lambda/3$ and raises degree by at
most four, giving omitted ground norm
$t_N=a^{\lfloor N/4\rfloor+1}/(1-a)$ and energy error
$\delta_N\le20\lambda t_N/(1-t_N)$.  The centered finite-window bound is
\begin{equation}
 \norm{(e^{-\sigma(L-\epsilon)}-e^{-\sigma(A_N-\mu_N)}P_N)P_{d_0}}
 \le\min\{2,2e^z z^r/r!+\sigma\delta_N\}.\qquad\text{\workstar}
 \label{recent:x-degree}
\end{equation}
For $d_0=0,N=20,0\le\sigma\le1$, it is below $10^{-6}$ at the cap, but that
larger matrix also remained unassembled.  The cutoff, input, and time sets in
\eqref{recent:x-delayed} and \eqref{recent:x-degree} are different.

Both constructions identify the same actual rank-21 sector at $R=9/2$ or
$N=4$: $\Omega=1$ and $\phi_p=2W_p$.  Gauge support completeness below this
energy follows from the four-cycle argument; five-edge supports cannot occur
on the bipartite graph.  Haar normalization and the absence of a zero parity
mask among triples of faces give the complete matrix
\begin{equation}
 (A_{21})_{00}=c,\quad (A_{21})_{0p}=-\lambda/2,\quad
 (A_{21})_{pq}=(3+c)\delta_{pq},\qquad
 \mu=c-\frac{\sqrt{9+20\lambda^2}-3}{2}.\qquad\text{\workstar}
 \label{recent:matrix21}
\end{equation}
There are nineteen dark face eigenvectors at $3+c$ and one bright excited
root $c+(3+\sqrt{9+20\lambda^2})/2$.  The spin-one character
$\chi_p=4W_p^2-1$ is omitted, normalized, and has energy 8, with
$\langle\chi_p,L\phi_p\rangle=-\lambda/2$.  Thus the small matrix is a
compression with genuine physical leakage.  X1's parent practical-truncation
gate remains limited.  Sources: \repo{research/round24/forward/x1/report.md},
\repo{research/round24/reverse/x1/report.md},
\repo{research/round24/skeptic/x1-review.md}.

\subsubsection{\workstar\ X2: complete residual, spectral data, and evaluated late heat}

Write $D=\sqrt{9+20\lambda^2}$, $w=(D-3)/2$,
$h=\lambda/[2(3+w)]$, $Z=1+20h^2$, and
$f=(\Omega+h\sum_p\phi_p)/\sqrt Z$.  The full residual is
\begin{equation}
 F=S^2-5,\qquad (L-\mu)f=-\frac{2\lambda h}{\sqrt Z}F,
 \qquad \rho^2=\frac{195\lambda^4}{4(3+w)D}.\qquad\text{\workstar}
 \label{recent:residual21}
\end{equation}
To obtain its complete norm, use
$\mathbb E W_p^4=1/8$ and
$\mathbb E W_p^2W_q^2=1/16$ for distinct faces, even when they share an edge:
integrating an exclusive link first makes each squared factor constant.
No four distinct faces have zero combined center parity.  Hence
\begin{equation}
 \mathbb E S^2=5,\quad \mathbb E S^4=295/4,\quad
 F=\tfrac14\sum_p\chi_p+2\sum_{p<q}W_pW_q,\quad
 \norm F^2=195/4.
 \label{recent:haar-residual}
\end{equation}
This accounts for 20 spin-one and 190 orthogonal pair channels, not the whole
omitted Hilbert space.  Keeping only spin-one channels loses a factor 39 in the
squared residual.  The forward-origin electric refinement resolves the 62
shared-edge products into singlet/triplet norms $1/64,3/64$, whereas 128
edge-disjoint products have energy 6.  Its complete weight table is
\begin{center}
\begin{tabular}{c|rrrr}
Electric energy & $9/2$ & $6$ & $13/2$ & $8$\\\hline
Squared weight in $F$ & $31/8$ & $32$ & $93/8$ & $5/4$
\end{tabular}
\end{center}
The weights sum to $195/4$ and their first moment is 295.  The independent
form identity $\langle F,KF\rangle=4\mathbb E S^4=295$ follows by integration
by parts from $KS=3S$.  A link-center transformation sends every face to its
negative, so the remaining odd potential moment vanishes.  Thus for
$\xi=-F/\sqrt{195/4}$, $\langle\xi,L\xi\rangle=\nu=236/39+c$.

The true ground $\epsilon$ is simple and every other eigenvalue is at least
3, by min-max.  Taking the spectral expectation of
$(L-\epsilon)(L-3)\ge0$ in $f$, and using the actual variational plane
$\operatorname{span}\{f,\xi\}$, gives
\begin{align}
 \delta_-:=\frac{2\rho^2}{\nu-\mu+\sqrt{(\nu-\mu)^2+4\rho^2}}
 &\le\mu-\epsilon\le\delta_+:=\frac{\rho^2}{3-\mu},\notag\\
 \norm{G-|f\rangle\langle f|}&\le p:=\frac\rho{3-\mu}.
 \qquad\text{\workstar}
 \label{recent:temple21}
\end{align}
Reverse obtains the valid weaker lower displacement
$\rho^2/(8+40\lambda-\mu+\rho)$ from complete kinetic support at energies at
most 8.  At $\lambda=0.01$ the forward enclosure is
$0.19983332325\le\epsilon\le0.19983333365$ and $p<8.32\cdot10^{-5}$.
Neither trial plane is asserted invariant.

Decomposing each heat operator into its own ground projection and excited
remainder yields
\begin{equation}
 \norm{e^{-\sigma(L-\epsilon)}-e^{-\sigma(A_{21}-\mu)}P_{21}}
 \le p+e^{-(3-\mu)\sigma}+e^{-(3+D)\sigma/2}.
 \label{recent:late21}
\end{equation}
A continuous-coupling envelope and certified scalar evaluator make this less
than $0.000085$ for every $\lambda\in[0,0.01]$, $\sigma\ge5$.  Forward's distinct
cap-only example begins at $\sigma=4$ with bound $0.000103$.
The actual finite representation is
$G_R+e^{-D\sigma}(U-G_R)+e^{-(3+D)\sigma/2}(P_{21}-U)$, where $U$ projects
onto the vacuum/symmetric-face plane.  The computed rational $\widehat G_R$ is
exactly idempotent and the ground coefficient stays exactly one.  Positive
gap rounding and scalar exponential errors are budgeted separately.  The
reverse scalar algorithm uses the degree-160 positive Taylor enclosure for
$e^x$ on $[0,96]$, reciprocates, and returns zero beyond 96.  Its uniform
$e^{-x}$ error is below $3\cdot10^{-30}$, since the reciprocal width is bounded
by $96(80!)^2/[161!(1-96/162)]$ and the omitted tail by $80!/96^{80}$.
The combined representation allowance is below $3\cdot10^{-25}$.
Sources: \repo{research/round24/forward/x2/report.md},
\repo{research/round24/reverse/x2/report.md},
\repo{research/round24/skeptic/x2-review.md}.

\subsubsection{\workstar\ Z1--Z2: a relative obstruction and a preparation-based repair}

Let $E=e^{-\sigma(L-\epsilon)}$ and
$E_R=e^{-\sigma(A_{21}-\mu)}P_{21}$.  The relative quantity is
$\norm{(E-E_R)x}/\norm{Ex}$ for the same nonzero retained input.  Heat injectivity
makes the denominator positive at every finite time, but not uniformly in time.
The exact vacuum row $L\Omega=c\Omega-\lambda S$ lies in $P_{21}$, and
\eqref{recent:temple21} proves $\delta=\mu-\epsilon>0$ at positive coupling.
For $u=(L-\epsilon)\Omega$ one has $Gu=0$ but
$\langle f,u\rangle=\delta\langle f,\Omega\rangle>0$, whence
\begin{equation}
 \frac{\norm{(E-E_R)u}}{\norm{Eu}}
 \ge\frac{\delta\langle f,\Omega\rangle}{\norm u}
 e^{(3-\mu)\sigma}-1\longrightarrow\infty.\qquad\text{\workstar}
 \label{recent:relative-obstruction}
\end{equation}
This is a fixed actual retained vector for each fixed $\lambda>0$, although
preparing its exact ground-orthogonal coefficient requires unknown $\epsilon$.
The computable $v=(A_{21}-\mu)\Omega$ is instead a bright Ritz eigenvector with
true ground overlap $-\delta\langle\psi,\Omega\rangle\ne0$; its true-relative
error tends to one, while its Ritz-output-relative error diverges.
At the cap forward certifies the first error above 6.9 at $\sigma=6$, and the
second within $10^{-4}$ of one at $\sigma=8$.  The mechanism is distinct ground
overlap, not divergence of either correctly centered heat operator.  At
$\lambda=0$ or $\sigma=0$ the retained physical error is exactly zero.
Sources: \repo{research/round24/forward/z1/report.md},
\repo{research/round24/reverse/z1/report.md},
\repo{research/round24/skeptic/z1-review.md}.

In the remaining heat estimates, integration coordinates $t,s$ and joining values $T$ are dimensionless heat coordinates, equal in units to $\sigma$; a duration $t$ corresponds to physical time $(\hbar/\alpha)t$.

Z2 reconstructs the entire retained-to-omitted map, not just its ground column.
With $\eta_{pq}=4W_pW_q$, the actual orthogonal residual channels give
\begin{align}
 B\Omega&=0,\quad
 B\sum_pb_p\phi_p=-\frac\lambda2\left[
 \sum_pb_p\chi_p+\sum_{p<q}(b_p+b_q)\eta_{pq}\right],\notag\\
 (B^*B)|_{\mathrm{faces}}&=\frac{\lambda^2}{4}(19I+J_{20}),
 \qquad\norm B=\frac{\sqrt{39}}2\lambda.\qquad\text{\workstar}
 \label{recent:omitted-gram21}
\end{align}
This follows by expanding the squared norm on arbitrary complex coefficients.
For the exact retained vacuum evolution,
\begin{equation}
 BE_R(s)\Omega=-\frac{\lambda^2}{D}(1-e^{-Ds})F,
 \qquad\norm{BE_R(s)\Omega}=\beta(1-e^{-Ds}),\quad
 \beta=\frac{\lambda^2\sqrt{195}}{2D}.
 \label{recent:vacuum-loading21}
\end{equation}
The zero initial omitted column therefore does not remain zero.  In fact
$Q(L-\epsilon)^2\Omega=\lambda^2F$.
The exact centered Duhamel defect is $B+\delta P$, not $B$ alone.  Thus
\begin{equation}
 V(t)=\beta[t-(1-e^{-Dt})/D]+\delta_+t
 \quad\text{bounds}\quad\norm{(E(t)-E_R(t))\Omega}.
 \label{recent:early21}
\end{equation}
For $f_0=\langle f,\Omega\rangle$ and
$q_0=\norm{(I-G_R)\Omega}$, an independent spectral bound is
$J(t)=p+(q_0+p)e^{-(3-\mu)t}+q_0e^{-Dt}$, and the true denominator is at least
$f_0-p$.  The increasing $V$ and decreasing $J$ can be joined:
\begin{equation}
 \sup_{t\ge0}\frac{\norm{(E(t)-E_R(t))\Omega}}{\norm{E(t)\Omega}}
 \le\frac{\max\{V(T),J(T)\}}{f_0-p}.\qquad\text{\workstar}
 \label{recent:join21}
\end{equation}
For normalized $x\in P_{21}$ with $\norm{x-\Omega}\le\eta$, add to the early
bound $\eta[(\rho+\delta_+)t+\norm B(1-e^{-3t})/3]$, and use late bound
$p+(q_0+\eta+p)e^{-(3-\mu)t}+(q_0+\eta)e^{-3t}$ with denominator
$f_0-\eta-p$.  This preserves the same prepared input in both evolutions.
The continuous envelopes
$\rho,\beta\le\sqrt{195}\Lambda^2/6$, $g\ge3-20\Lambda$,
$q_0\le\sqrt5\Lambda/3$, and $f_0\ge\sqrt{1-5\Lambda^2/9}$ cover every
$0\le\lambda\le\Lambda\le0.01$, rather than a parameter grid.  They give
\begin{center}
\begin{tabular}{lcc}
Input class & joining $\sigma$ & all-time relative bound\\\hline
Vacuum & $3/2$ & $0.00028$\\
Normalized $P_{21}$ ball, radius $0.01$ & $8/5$ & $0.00044$\\
Normalized retained Ritz overlap $\ge0.99$ & $2$ & $0.002$
\end{tabular}
\end{center}
Reverse's ball without a normalization requirement has the distinct bound
$0.000720$ at join $29/10$; its vacuum bound is $0.000600$.  These are not
averaged with the sharper normalized results.  Ritz overlap $a>p$ provides the
checkable denominator floor $a-p$; uncertainty in overlap, coefficients, or
preparation must be subtracted first.  Comparing to an ideal vacuum instead
of the same prepared input adds the preparation bias separately.  Rational
projectors and positive-gap scalar evaluation add negligible uniform numerical
cost while preserving an exact ground coefficient.  Z1's excluded ground-null
inputs remain excluded.
Sources: \repo{research/round24/forward/z2/report.md},
\repo{research/round24/reverse/z2/report.md},
\repo{research/round24/skeptic/z2-review.md}.

\subsubsection{\workstar\ AC1--AC2: a 293-state physical space and delayed new-channel loading}

Adjoin every old residual channel to $P_{21}$ and close under $K$.  A shared-edge
pair must be split into its spin-zero and spin-one common-edge parts.  In exact
rational coordinates the basis and metric are
\begin{center}
\begin{tabular}{lrrr}
Vector type & count & norm squared & electric energy\\\hline
$\Omega$ &1&1&0\\
$\phi_p=2W_p$&20&1&3\\
$\chi_p=4W_p^2-1$&20&1&8\\
$d_{pq}=4W_pW_q$, edge-disjoint&128&1&6\\
$s_{pq}=8X_{pq,s}$, shared-edge singlet&62&1&$9/2$\\
$t_{pq}=8X_{pq,t}$, shared-edge triplet&62&3&$13/2$
\end{tabular}
\end{center}
Thus $P_{293}$ has rank $1+20+20+128+124=293$.  Orthogonality follows from
pair center parity, distinct spin-one supports, and orthogonal common-edge
Casimir projections.  This is a generated $K$-reducing subspace, not an entire
electric shell; omitted intertwiners remain in $Q_{293}$.

The forward-origin complete magnetic matrix follows from center parity and the
exact old columns
\begin{align}
 S\Omega&=\tfrac12\sum_p\phi_p,\notag\\
 S\phi_p&=\tfrac12\Omega+\tfrac12\chi_p
 +\tfrac12\sum_{q\text{ edge-disjoint}}d_{pq}
 +\tfrac14\sum_{q\text{ sharing}}(s_{pq}+t_{pq}).\qquad\text{\workstar}
 \label{recent:matrix293}
\end{align}
All new vectors are even under simultaneous face-sign reversal and all old
face vectors are odd, so same-parity magnetic entries vanish.  Metric
self-adjointness fixes the reverse entries, notably $3/4$ from a triplet to
each face, not $1/4$.  The assembled matrix has 1088 directed nonzero magnetic
entries and diagonal Gram with 231 weights one and 62 weights three.
Its complete omitted map satisfies
\begin{equation}
 B_+=Q_{293}LP_{293}=-\lambda Q_{293}SP_{293},\quad
 B_+P_{21}=0,\quad\norm{B_+}\le20\lambda.
 \label{recent:new-leakage}
\end{equation}
The bound covers all 272 new columns together, including unexpanded omitted
intertwiners; it is not an evaluation of their full Gram.

Because the old Ritz residual is now retained, its angle to the new Ritz
ground is at most $p_{21}=\rho_{21}/(3-\mu_{21})$.  Since $B_+f_{21}=0$,
\begin{equation}
 \rho_+\le20\lambda p_{21},\qquad
 \norm{G-G_+}\le\frac{\rho_+}{3-\mu_+},\qquad
 0\le\mu_+-\epsilon\le\frac{\rho_+^2}{3-\mu_+}.
 \label{recent:ground293}
\end{equation}
At the coupling cap, conservative exact envelopes are
$p_{21}\le1/12000$, $\rho_+\le1/60000$,
$p_+\le1/168000$, and $\delta_+\le1/10080000000$.
AC1 improves ground certification but its coarse all-time ball bound $0.0014$
is worse than Z2's $0.00044$; increasing dimension alone does not guarantee a
better error certificate.  Its parent accuracy gate was therefore limited.
Sources: \repo{research/round26/forward/ac1/report.md},
\repo{research/round26/reverse/ac1/report.md},
\repo{research/round26/skeptic/ac1.md}.

AC2 exploits the actual loading sequence.  Let $N=P_{293}-P_{21}$,
$P_f=P_{21}-|\Omega\rangle\langle\Omega|$, $A_+=P_{293}LP_{293}$, and
$C=NA_+P_{21}$.  Exact matrix identities give
$NA_+N=D_N+cN$, $D_N\ge9/2$, $C\Omega=0$, and
$\norm C=\sqrt{39}\lambda/2$ from \eqref{recent:omitted-gram21}.
For original inputs $x\in P_{21}$, the new retained component
$y(t)=Ne^{-t(A_+-\mu_+)}x$ obeys the full equation
\begin{equation}
 y'=-(D_N+c-\mu_+)y-CP_fe^{-t(A_+-\mu_+)}x,\qquad y(0)=0.
 \label{recent:loading-equation}
\end{equation}
Its forcing uses the entire retained trajectory and therefore includes feedback.
For $\norm x=1$, $\norm{x-\Omega}\le\eta=0.01$ and
$0\le\lambda\le\Lambda\le0.01$, put
\begin{equation}
 g=3-20\Lambda,\quad p_{21}=\frac{7\Lambda^2}{3g},\quad
 q_+=\frac{3\Lambda}{4}+\frac{3p_{21}}2,\quad b=q_++\eta,
 \quad\overline C=\frac{25\Lambda}{8},\quad d=9/2.
 \label{recent:loading-envelopes}
\end{equation}
The new ground/excited decomposition bounds the forcing by
$q_++be^{-gt}$.  Convolution yields
\begin{equation}
 \norm{y(t)}\le\overline C\left[
 q_+\frac{1-e^{-dt}}d+b\frac{e^{-gt}-e^{-dt}}{d-g}\right].
 \label{recent:delayed-loading}
\end{equation}
Duhamel with the complete $B_+$ and its separate centering defect consequently
has increasing envelope
$\overline V(t)=\delta_+t+20\Lambda\overline C[q_+t/d+b/(gd)]$.
The independent late envelope is
$J(t)=p_++(2b+p_+)e^{-gt}$.  The true full output is bounded below using the old
full-ground result,
$D_\eta=1-5\Lambda^2/9-\eta-p_{21}>0$.
At $\Lambda=0.01$, the join $t_0=13/5$ and $e^{-182/25}<1/1400$ give
\begin{equation}
 \begin{aligned}
 D_\eta&=\frac{7127}{7200},\qquad
 \overline V(t_0)=\frac{457097}{12600000000},\qquad
 J(t_0)\le\frac{2441}{78400000},\\
 \sup_{\sigma\ge0}\frac{\norm{(E-E_+)x}}{\norm{Ex}}
 &\le\frac{457097}{12472250000}<0.000037.\qquad\text{\workstar}
 \end{aligned}
 \label{recent:ac-relative}
\end{equation}
All envelopes cover the continuous coupling interval.  Reverse independently
uses face-block norms $\norm{NSP_f}=\sqrt{39}/2$ and
$\norm{P_fS(I-P_f)P_{293}}=\sqrt{59}/2$ to obtain a delayed-leakage integral
\begin{equation}
 \int_0^t\norm{B_+E_+(s)x}\,ds\le\frac{130}{9}\lambda^2
 \left[\eta\frac{1-e^{-3t}}3+\frac{4\lambda}{3}
 \left(t-\frac{1-e^{-3t}}3\right)\right],
 \label{recent:reverse-loading}
\end{equation}
leading to $0.000048$ for the ball and $0.000043$ for the vacuum.
The sharper result is confined to original $P_{21}$ inputs: new retained
preparations would not have $y(0)=0$.  At $\sigma=0$ and at $\lambda=0$ the
physical retained error is exactly zero.  AC2 certifies the exact retained
semigroup; numerical evaluation was still a separate task.
Sources: \repo{research/round26/forward/ac2/report.md},
\repo{research/round26/reverse/ac2/report.md},
\repo{research/round26/skeptic/ac2.md}.

\subsubsection{\workstar\ AF1--AF2: execute the heat and prevent secular centering error}

AF1 fixes $\lambda=0.01$, $\sigma=1$, and input $\Omega$, and exports all 293
coordinates.  Every norm uses the actual diagonal Gram $\Gamma$, with
$\Tr\Gamma=417$.  A trial residual identifies the minimum only after the
spectral separation $\lambda_2(A_+)\ge3$ and $\mu_+\le1/5$ are proved.
For a rational trial vector $v$, define its Rayleigh quotient $a$ and residual
$\rho_v^2=\norm{A_+v}_\Gamma^2/\norm v_\Gamma^2-a^2$.  Then
\begin{equation}
 a-\frac{\rho_v^2}{3-a}\le\mu_+\le a,\qquad
 \norm{P_v-G_+}\le\frac{\rho_v}{3-a},\quad
 P_vx=v\frac{\langle v,x\rangle_\Gamma}{\langle v,v\rangle_\Gamma}.
 \label{recent:numerical-minimum}
\end{equation}
The first inequality is again the full retained spectral-measure argument,
not a small residual near an unidentified eigenvalue.

Forward uses an explicit one-step electric residual correction, a center radius
below $10^{-11}$, and the degree-100 polynomial in $A_+-\widehat\mu$;
reverse uses the old Ritz energy interval, center radius below $10^{-8}$,
and degree 80.  Exact integer or rational Horner arithmetic has no accumulated
floating roundoff.  Since $\norm{A_+-\widehat\mu}\le9$, their respective
polynomial bounds are $2^{14}9^{101}/101!$ and $2\cdot9^{81}/81!$.
Coordinate grids $10^{-50}$ and $10^{-30}$ cost at most
$21/(2\cdot10^{50})$ and $11\cdot10^{-30}$ in the physical metric.  The separate
center cost is $d_\mu/(1-d_\mu)$ at this time.
Adding the independently retained physical omission and dividing by the true
floor $7199/7200$ gives full relative point bounds $0.000015$ forward and
$0.000014$ reverse.  Exact Gram comparisons put their exported-vector difference
inside the sum of their independent numerical radii.
Sources: \repo{research/round26/forward/af1/report.md},
\repo{research/round26/reverse/af1/report.md},
\repo{research/round26/skeptic/af1.md}.

AF2 prevents the uncertain center from being exponentiated indefinitely.
It fixes the same coupling and original normalized radius-0.01 $P_{21}$ class,
and defines a piecewise approximation:
\begin{equation}
 \widehat E_+(\sigma)x=
 \begin{cases}
  P_n(\sigma(A_+-\widehat\mu))x,&0\le\sigma\le8,\\
  P_vx,&\sigma>8,
 \end{cases}
 \qquad P_n(H)=\sum_{j=0}^n\frac{(-H)^j}{j!}.
 \label{recent:piecewise-solver}
\end{equation}
Forward uses degree 256 with uniform Taylor error
$2\cdot72^{257}/257!<10^{-25}$, center cost
$8d_\mu/(1-8d_\mu)$, and complex export allowance $15\cdot10^{-40}$.
Its corrected rational trial is certified by \eqref{recent:numerical-minimum},
independently of the heuristic correction's convergence.  Reverse constructs
$v=(I-A_+/9)^{64}\Omega$ exactly.  The excited-to-ground power ratio is at most
$15/22$ and vacuum ground overlap squared is at least $14/15$; the final
Rayleigh residual certifies the actual resulting vector.  It uses degree 240,
Taylor bound $16\cdot72^{241}/241!$, and real export allowance
$11\cdot10^{-30}$.

The late branch has no uncertain ground exponent.  For every $\sigma>8$,
\begin{equation}
 \norm{e^{-\sigma(A_+-\mu_+)}-P_v}
 \le\frac{\rho_v}{3-a}+e^{-(14/5)\sigma}.
 \label{recent:late-solver}
\end{equation}
The full retained excited tail is below $1/(5\cdot10^9)$ in forward's rational
certificate and below $10^{-9}$ in reverse's.  Each complete numerical
allowance $e_{\rm num}$ is below $10^{-8}$.  Therefore
\begin{equation}
 \frac{\norm{\widehat E_+(\sigma)x-E(\sigma)x}}{\norm{E(\sigma)x}}
 \le\frac{457097}{12472250000}
   +\frac{e_{\rm num}}{7127/7200}<0.000037
 \quad(\sigma\ge0).\qquad\text{\workstar}
 \label{recent:computed-relative}
\end{equation}
The two terms are physical compression and retained computation errors.
The sample exports include the exact normalized preparations
$(159999\Omega+800\phi_0)/160001$ and, forward only,
$(159999\Omega+800i\phi_{19})/160001$, whose squared vacuum distance is
$4/160001<10^{-4}$.  Forward exports eighteen full complex vectors at
$\sigma=0,1,13/5,8,9,1000$; reverse exports twelve real vectors, replacing the
last time by $10^6$.  Continuous coverage comes from the two analytic branch
bounds, not interpolation between these exports.  Outputs are not renormalized.
The reverse command interface accepts exact rational times and 293 real rational
coordinates, enforcing original support, normalization, and radius; the operator
estimate covers complex vectors, but that does not make the real interface a
complex interface.  Fixed-coupling computation is not continuous-coupling
execution.  Time, coefficient, or preparation measurement uncertainty requires
its own budget.

The remaining finite-graph tasks are growing-graph control and additional
observable/input classes with complete new omission certificates.  Neither the
all-time heat result nor its small error implies real-time accuracy, a
homogeneous numerical gap, physical scale matching, or a continuum construction.
Sources: \repo{research/round26/forward/af2/report.md},
\repo{research/round26/reverse/af2/report.md},
\repo{research/round26/skeptic/af2.md}.

\subsection{Canonical endpoint observables and an exact reached spectrum}
\label{recent:canonical}

Use the canonical summable full-link model, with complete selected-strip and
free-link reference factors, unique reference ground $\Omega$, and
\begin{align}
 H_q&=H_0+V_q,\qquad
 V_q=-\frac{\alpha\tau_q}{24}\sum_{f\text{ omitted}}q^{|a_f|_1}x_f,
 \quad x_f=\tfrac12\Tr U_f,\notag\\
 b(q)&=\frac{p(q)}{24(1-q)^3(1+q)^2(1+q^2)},\qquad
 p(q)=2+5q+5q^2+6q^3+3q^4,\notag\\
 \tau_q&=\frac{\eta}{8b(q)},\qquad0<q,\eta<1.
 \label{recent:canonical-profile}
\end{align}
The perturbation is norm convergent with $\norm{V_q}=\alpha\eta/8$ and
$D(H_q)=D(H_0)$.  The admitted stationary ground projection satisfies
\begin{equation}
 d_q=\norm{P_q-P_0}\le\widehat d_q:=\frac{\sigma_q}{\bar g},\quad
 \sigma_q^2=\frac{\alpha^2\tau_q^2b(q^2)}{96},\quad
 \bar g=\frac{\alpha(1-\eta)}8,
 \qquad\widehat d_q=O((1-q)^{3/2}).
 \label{recent:canonical-state}
\end{equation}
For the convention $\beta_q^t(B)=e^{-itH_q/\hbar}Be^{itH_q/\hbar}$, define
$C_q^B(t)=\omega_q(B\beta_q^t(B))-\omega_q(B)^2$.
All endpoint conclusions in this subsection are restricted to
\begin{equation}
 T_q=\frac z\eta\frac\hbar\alpha(1-q)^{-3},\qquad
 0<z\le10^{-6},\qquad
 s_q=\alpha\tau_qT_q/\hbar\longrightarrow\frac{8z}{7},\quad s_q\le\frac{3z}{2}.
 \label{recent:endpoint-clock}
\end{equation}
The variable $s_q$ is a proof coordinate at the original physical time $T_q$.
The trigonometric formulas below have algebraic extensions, but no full
stationary endpoint theorem beyond the displayed $z$ interval is asserted.

\subsubsection{\workstar\ U1--U2: a first-order resonance becomes a full-system witness}

Let $O=(3,1,0)$ and $D=(3,2,1)$, and take paths
$X:O\to(3,2,0)\to D$, $Y:O\to(3,1,1)\to D$, and
$Z:O\to(4,1,0)\to(4,2,0)\to(4,2,1)\to D$.
Their eight links are free singleton factors.  The two real characters
$\phi_X=\Tr(XZ^{-1})$, $\phi_Y=\Tr(YZ^{-1})$ are orthonormal and have energy
$E=6(3\alpha/4)=9\alpha/2$.  Set
$\psi_+=(\phi_X+\phi_Y)/\sqrt2$ and
$W=(|\psi_+\rangle\langle\Omega_F|+\mathrm{adjoint})\otimes I_{F^c}$.
It is a fixed gauge-invariant bounded local operator, with $W^3=W$,
$\norm W=1$, and reference variance one.  It is not a Wilson multiplier.

The yz face at $O$ has $x_{f_0}=\Tr(XY^{-1})/2$.
Haar character convolution gives
$\langle\phi_Y,x_{f_0}\phi_X\rangle=1/4$.  Every other face has an unmatched
free-link center parity, and all diagonal terms vanish by the same selection
rule.  The actual full-model matrix element is therefore
\begin{equation}
 \langle\psi_+,V_q\psi_+\rangle=-\frac{\alpha\tau_qq^4}{96}.
 \qquad\text{\workstar}
 \label{recent:u-resonance}
\end{equation}
It gives first-order phase-removed response $is_qq^4/96\to iz/84$.
The vacuum subtraction uses $\langle\Omega,V_q\Omega\rangle=0$, not
$V_q\Omega=0$.  A derivative alone did not establish an endpoint; U1's gate is
limited and its two-state plane is not assumed invariant.

U2 supplies the complete remainder.  Each complete reference factor owns at
most ten links and meets at most forty omitted faces; one face reaches at most
four factors.  Starting from eight factors, the ordered connected word count
at depth $n$ is bounded by $\prod_{k=0}^{n-1}40(8+3k)$.  Combining coupling
$\alpha\tau_q/24$, commutator factor two and simplex $|t|^n/n!$ gives
\begin{equation}
 \norm{\text{order }n}\le\frac{(8/3)_n}{n!}
 \left(10\alpha\tau_q|t|/\hbar\right)^n.
 \label{recent:u-connected}
\end{equation}
Strong bounded-generator Dyson evolution first justifies the expansion for each
$q<1$; connected counting then supplies its useful uniform bound without
assuming operator-norm continuity through the unbounded $H_0$.
For $x=10\alpha\tau_q|t|/\hbar<1$,
\begin{equation}
 R(x)=(1-x)^{-8/3}-1-\tfrac83x
 \le\widehat R(x):=\frac{44x^2}{9(1-x)^5}.
 \label{recent:u-tail}
\end{equation}
Taylor's integral remainder proves this using the second derivative
$(88/9)(1-x)^{-14/3}$; reverse's independent algebraic bound within the same
author's reconstruction is $x^2(6-8x+3x^2)/(1-x)^3$.
State replacement in the bounded connected expression costs $6\widehat d_q$,
so
\begin{align}
 \left|e^{iEt/\hbar}C_q^W(t)-1-is_qq^4/96\right|
 &\le6\widehat d_q+\widehat R(10s_q),\notag\\
 \liminf_{q\uparrow1}|C_q^W(T_q)-C_0^W(T_q)|
 &\ge\frac z{84}-\frac{44}{9}
 \frac{(80z/7)^2}{(1-80z/7)^5}>\frac z{168}>0.
 \qquad\text{\workstar}
 \label{recent:u-endpoint}
\end{align}
The final strict inequality follows by monotonicity of $\widehat R(80z/7)/z$
and exact arithmetic at $z=10^{-6}$.  This makes the earlier sufficient
$\gamma<3$ time-window exponent sharp in the existence-of-one-observable sense;
it does not decide every observable or compute the limiting trajectory.
Both U loops use same-author self-review.
Sources: \repo{research/round23/forward/u1/report.md},
\repo{research/round23/reverse/u1/report.md},
\repo{research/round23/forward/u2/report.md},
\repo{research/round23/reverse/u2/report.md},
\repo{research/round23/skeptic/u1.md},
\repo{research/round23/skeptic/u2.md}.

\subsubsection{\workstar\ V1--V2: finite outcomes are not a finite implementation}

The exact spectral projections
$P_\pm=(W^2\pm W)/2$, $P_0=I-W^2$ give outcomes $-1,0,+1$, with an
infinite-dimensional local zero eigenspace.  For self-adjoint
$\norm{A-W}\le\delta$, $\norm A\le1+\delta$, two applications of the product
difference inequality, including the connected means, prove
\begin{equation}
 |C_q^A-C_q^W|\le4\delta+2\delta^2,\qquad
 |(C_q^A-C_0^A)-(C_q^W-C_0^W)|\le8\delta+4\delta^2.
 \label{recent:readout-transfer}
\end{equation}
This is uniform in time.  Choosing fixed $\delta=z/100000$ retains endpoint
margin greater than $z/200$; for contractions the two-model cost improves to
$8\delta$.  A binary first-moment POVM instead has outcome-square one and does
not identify $\omega(W^2)$ with its outcome variance.

Ordinary sequential projective measurements cannot generally recover the
complex unordered correlation.  Set the self-adjoint unitary completions
$R_\pm=W\pm P_0$.  Then
\begin{equation}
 \omega(W\beta^t(W))=\frac14\sum_{s,r\in\{+,-\}}
 \omega(R_s\beta^t(R_r)).\qquad\text{\workstar}
 \label{recent:coherent-readout}
\end{equation}
Four controlled-unitary products and two ancilla quadratures require eight
binary settings, plus a separate mean measurement.  This applies the established
interferometric primitive; it assumes coherent access to the actual full
both-sign evolution and the intended stationary states.  The exact PVM and
instrument identity do not construct those resources.

A separate V1 obstruction concerns the full local Haar operator norm.
For every multiplication operator $M_g$, a weakly-null orthonormal sequence
in a positive-measure level set of $|g|$ gives
$\norm{W-M_g}\ge\norm g_\infty$ because $W_F$ has finite rank.  The triangle
inequality also gives $\norm{W-M_g}\ge1-\norm g_\infty$, hence
\begin{equation}
 \inf_g\norm{W_F-M_g}\ge\frac12.\qquad\text{\workstar}
 \label{recent:multiplier-obstruction}
\end{equation}
The gate admits this in the full local Haar norm, not an unconstructed
physical-quotient norm.  It blocks this norm-transfer route to Wilson observables,
not the existence of a different Wilson witness.

For V2, the actual eigenvectors defining $W$ prove domain preservation and
\begin{equation}
 \norm{[H_0,W]}=9\alpha/2,\qquad
 \norm{[H_q,W]}\le\alpha(9/2+\eta/4).
 \label{recent:clock-commutator}
\end{equation}
Graph closure extends the finite-factor commutator, and strong integration
then gives a norm-Lipschitz physical-time bound for this particular observable.
For one predeclared endpoint, eighteen independent bounded-output settings,
$N$ trials per setting, and tolerance $r$, the exponential-moment argument and
union bound give failure probability at most $36e^{-Nr^2/2}$.
In forward's convention, full trace-norm preparation error $s$, complete-setting
operator error $g$, and a common per-model time offset
$|\Delta t|\le(\hbar/\alpha)\theta$ yield
\begin{equation}
 N\ge\frac2{r^2}\log\frac{36}{\kappa},\qquad
 |\widehat D-(C_q^W(T_q)-C_0^W(T_q))|
 \le8(r+s+g)+(9+\eta/4)\theta
 \label{recent:readout-budget}
\end{equation}
with confidence at least $1-\kappa$.  Reverse permits per-shot timing jitter,
uses half trace distance $\varepsilon_p$, and obtains instead
$8(r+2\varepsilon_p+\varepsilon_i)
+2(9+\eta/4)\alpha h/\hbar$; the timing coefficients are not interchangeable.

At $\eta=1/2,z=10^{-6},q=1-10^{-6}$, forward takes
$r=s=g=\theta=10^{-10}$ and $N=1.8\cdot10^{21}$, giving $3.24\cdot10^{22}$
preparations, a finite-$q$ true margin above $10^{-8}$, and confidence radius
$3.3125\cdot10^{-9}$ at physical time $2\cdot10^{12}\hbar/\alpha$.
Reverse's different example, $q=1-10^{-8}$, has radius $8.4225\cdot10^{-10}$
and time $2\cdot10^{18}\hbar/\alpha$.  Neither is an implemented or efficient
experiment, an optimal resource lower bound, or confidence over a continuum of
times.  Sources: \repo{research/round24/forward/v1/report.md},
\repo{research/round24/reverse/v1/report.md},
\repo{research/round24/skeptic/v1-review.md},
\repo{research/round24/forward/v2/report.md},
\repo{research/round24/reverse/v2/report.md},
\repo{research/round24/skeptic/v2-review.md}.

\subsubsection{\workstar\ Y1--Y2: spatial collars and an actual scalar limit}

Starting from the eight free factors, recursively include every complete owner
of every omitted face meeting the current factors.  Complete selected strips
must be included because reference evolution spreads an operator within its
whole factor.  Forward retains all omitted faces internally supported on the
new factor set; reverse retains those meeting the preceding set.  The conventions
therefore agree on factors and links but not on their finite Hamiltonians:
\begin{center}
\begin{tabular}{crrrr}
Depth & factors & links & forward faces & reverse faces\\\hline
1&35&98&20&19\\
2&135&297&129&125
\end{tabular}
\end{center}
Both are unrenormalized subsets of the actual $q$ profile, retain all Gauss
actions, and still have infinite-dimensional onsite Haar spaces.
Every connected word of length at most $k$ lies in collar $k$; longer retained
words cancel exactly when comparing the two expansions.  One full tail, not
two, therefore suffices:
\begin{equation}
 \sup_{|t|\le T_q}\norm{\beta_q^t(W)-\beta_{q,k}^t(W)}
 \le D_k(15z),\qquad
 D_k(x)=\frac{(8/3)_{k+1}x^{k+1}}{(k+1)!(1-x)^{k+4}}.
 \qquad\text{\workstar}
 \label{recent:spatial-tail}
\end{equation}
The bound includes every boundary excursion and return.  Reverse independently
uses the coefficientwise larger $(3)_n/n!$ and exact rational tail
\begin{equation}
 B_k(x)=\frac{x^{k+1}[(k+2)(k+3)-2(k+1)(k+3)x+(k+1)(k+2)x^2]}
 {2(1-x)^3}.
 \label{recent:reverse-spatial-tail}
\end{equation}
Each regional Hamiltonian has its own unique stationary ground.  Comparing both
states and both connected means costs
$6(d_q+d_{q,k})$, bounded by $12\widehat d_q$ in forward's route.
Thus the forward 98-link region retains limiting difference greater than $z/250$;
at $z=10^{-6}$ its spatial error is below $1.101\cdot10^{-9}$ and depth two
below $2.568\cdot10^{-14}$.  Finite spatial support does not make these evaluated
finite matrices.
Sources: \repo{research/round24/forward/y1/report.md},
\repo{research/round24/reverse/y1/report.md},
\repo{research/round24/skeptic/y1-review.md}.

For a fixed collar, its reference $K_k=H_{0,k}/\alpha$ is a finite sum of link
Casimirs plus bounded strip potentials and scalar ground shifts.  Peter--Weyl
projections below finite energy have finite rank; hence $K_k$ has compact
resolvent even before gauge reduction.  Let
$A_k=-\sum_{f\in O_k}x_f/24$ and retain the entire degenerate energy projections
in $\overline A_k=\sum_EP_EA_kP_E$.  This is a bounded self-adjoint strong direct
sum preserving $D(K_k)$; its vacuum block vanishes.

The needed averaging argument can be proved directly.  For a self-adjoint
pure-point $K$, bounded $A$, and $\ell>0$, put
$Z_\ell(s)=e^{isK/\ell}e^{-is(K/\ell+A)}$ and
$R_\ell(s)=\int_0^s(e^{iuK/\ell}Ae^{-iuK/\ell}-\overline A)du$.
On a fixed energy block the nonresonant coefficient is
\begin{equation}
 P_FR_\ell(s)P_E=
 \frac{\ell(e^{is(F-E)/\ell}-1)}{i(F-E)}P_FAP_E\quad(F\ne E).
 \label{recent:averaging-primitive}
\end{equation}
Orthogonal output energies and density imply strong convergence of this
primitive, uniformly on compact $s$ intervals and compact input sets.
Integration by parts in Duhamel makes it act only on the compact averaged orbit
and its bounded derivative, giving strong convergence of $Z_\ell$ and its
adjoint to $e^{-is\overline A}$.  If
$\Delta_E=\operatorname{dist}(E,\operatorname{spec}K\setminus\{E\})>0$,
the same proof yields the fixed-block estimate
\begin{equation}
 \sup_{|s|\le S}\norm{(Z_\ell(s)-e^{-is\overline A})P_E}
 \le\frac{2\ell M}{\Delta_E}(1+2SM),\qquad\norm A\le M.
 \label{recent:averaging-block}
\end{equation}
This is not a bound over all excited energies.  Compact-resolvent doublets with
vanishing splittings provide explicit counterexamples to a general
operator-norm averaging inference.

In the actual collar, $M_k=N_k/24$ and
$\norm{A_{q,k}-A_k}\le(1-q)\sum_{f\in O_k}|a_f|_1/24$.
At $\ell=\tau_q$ and the original $T_q$, the exact ordered scalar is
\begin{equation}
 e^{iEt/\hbar}\langle\Omega,W\beta_{q,k}^t(W)\Omega\rangle
 =\langle\psi_+,Z_{q,k}WZ_{q,k}^*\Omega\rangle.
 \label{recent:ordered-scalar}
\end{equation}
Both strong propagator limits are needed.  The averaged vacuum is fixed, so
for the entire energy-$9/2$ block
$B_k=P_{k,9/2}A_kP_{k,9/2}$ its regional limit is
$f_k(z)=\langle\psi_k,e^{-i(8z/7)B_k}\psi_k\rangle$.
The uniform spatial bound makes $f_k$ Cauchy, and choosing $k$ first and then
$q\uparrow1$ proves
\begin{equation}
 e^{iET_q/\hbar}C_q^W(T_q)\longrightarrow F_\infty(z),\qquad
 F_\infty(z)=\lim_{k\to\infty}f_k(z),\quad
 |F_\infty-f_k|\le D_k(15z).\qquad\text{\workstar}
 \label{recent:scalar-limit}
\end{equation}
No global $q=1$ generator or compact resolvent of the infinite reference is
used.  At this gate the actual matrices and their excited separations remained
unevaluated.  The forward-origin disk
\begin{equation}
 |F_\infty(z)-1-iz/84|
 \le D_2(15z)+\tfrac12(8z/7)^2(129/24)^2
 \label{recent:y-disk}
\end{equation}
has radius below $1.89\cdot10^{-11}$ at $z=10^{-6}$, by spectral Taylor
remainder and the known first moment $-1/96$.  It encloses an actual scalar,
not a computed trajectory.  This proof exceeds the finite-spectral bounded
reference hypotheses of the related norm-averaging theorem, so its pure-point
strong argument is essential.
Sources: \repo{research/round24/forward/y2/report.md},
\repo{research/round24/reverse/y2/report.md},
\repo{research/round24/skeptic/y2-review.md}.

\subsubsection{\workstar\ AA1--AA2: compute a reducing component, not the entire shell}

The unit cube translated by $O=(3,1,0)$ has twelve free links.  On its physical
energy-$9/2$ shell, quarter-energy arithmetic gives
$3n_{1/2}+8n_1+15n_{3/2}=18$.  Besides six half-spin links, the only solution
uses one half-spin and one spin-$3/2$ edge and is excluded by a gauge leaf.
At each cube vertex zero or two half-spin edges are required, so six occupied
edges form one simple six-cycle with a unique bivalent invariant intertwiner.
The sixteen six-cycles therefore form a complete orthonormal physical cube
shell.  Orthogonality follows from an unmatched edge-center sign.

A face can connect two such cycles only when their symmetric difference is its
four-edge boundary, replacing two adjacent edges by the other two.  The same
Haar convolution as U1 makes every allowed matrix element $1/4$; thus the
compressed potential is $-\mathsf A/96$, with $\mathsf A$ the exact sixteen-node
flip adjacency.  Enumeration and exact polynomial multiplication give
\begin{equation}
 \begin{aligned}
 \mathsf A(\mathsf A^2-4I)(\mathsf A^2-12I)&=0,\\
 \operatorname{spec}\mathsf A&=\{0^{(8)},2^{(3)},(-2)^{(3)},
  (2\sqrt3)^{(1)},(-2\sqrt3)^{(1)}\}.
 \qquad\text{\workstar}
 \end{aligned}
 \label{recent:aa-spectrum}
\end{equation}
Twelve vertices have degree two and four degree six.  The original coherent
source, reconstructed from its paths, has adjacency moments
$1,1,4,8,32,80,320$ through order six.  Polynomial spectral projectors give
weights $1/3$ at zero, $3/8$ and $1/8$ at $+2,-2$, and
$1/12\pm1/(8\sqrt3)$ at $\pm2\sqrt3$.  AA1 correctly remained limited because
cube-shell completeness alone did not exclude resonant exterior loading.

AA2 proves that missing inclusion.  A face external to the cube but sharing
cube edges shares exactly one, with twenty such actual faces.  Its opposite
edge is free and the remaining two sides have distinct exterior reference
owners.  The minimum $-3/4$ change of an occupied common edge is canceled by
$+3/4$ on the opposite free edge.  A remaining free side costs $3/4$; a strip
side creates an endpoint-charged vector orthogonal to its gauge-invariant
strip ground and costs at least $1/8$.  Distinct ownership therefore gives a
spectral support increment at least $1/4$, not merely a mean-energy inequality.
Cube-disjoint faces excite the exterior by at least $1/8$.  For internal faces
with $r$ occupied edges and $m$ spin-one fusion branches, the change is
\begin{equation}
 \Delta E=3-\tfrac32r+2m,
 \label{recent:aa-selection}
\end{equation}
whose nonzero values have magnitude at least $1/2$.
Consequently the sixteen-state space $\mathcal L$ times the unchanged exterior
ground reduces every sufficiently large regional pinched operator.  It need
not exhaust the whole energy shell, and the instantaneous interaction still
leaks.  All participating nonzero frequencies from $\mathcal L$ or the vacuum
are at least $1/8$ in magnitude, giving the evaluated primitive bound
$16\tau_qM_k$ and averaging error
$b_k=16\tau_qM_k(1+3zM_k)$.

These facts promote the AA1 cube scalar to the actual full canonical endpoint:
\begin{align}
 F_\infty(z)&=\frac13+\frac12\cos(2\theta)
 +\frac16\cos(2\sqrt3\theta)
 +i\left[\frac14\sin(2\theta)
 +\frac1{4\sqrt3}\sin(2\sqrt3\theta)\right],\notag\\
 \theta&=z/84,\qquad0<z\le10^{-6},\qquad
 |F_\infty-1-i\theta|\le2\theta^2=\frac{z^2}{3528}.
 \qquad\text{\workstar}
 \label{recent:aa-endpoint}
\end{align}
The complete second moment is four; it also gives
$|F_\infty-e^{i\theta}|\le3\theta^2/2$ by centering the spectral variable at
its mean one.  A two-state fit would lose the actual second and higher moments.

For finite $q$, keep the exact weighted flip matrix $Q(q)$, whose entries are
$q^4$ or $q^5$ and whose symmetric row norm is at most six.  Combining direct
full stationary-state replacement, the collar tail, and the two averaging
columns gives
\begin{equation}
 \left|e^{iET_q/\hbar}C_q^W(T_q)
 -\langle\psi_+,e^{is_qQ(q)/96}\psi_+\rangle\right|
 \le6\widehat d_q+D_k(15z)+32\tau_qM_k(1+3zM_k).
 \label{recent:aa-finite-q}
\end{equation}
A degree-eight rational scalar Taylor center adds at most
$(s_q/16)^9/9!$; comparison directly to the limiting formula separately adds
$s_q(1-q^5)/16+|s_q-8z/7|/16$.
At $\eta=1/2,z=10^{-6},q=1-10^{-12}$ and depth three with $N_3=332$, the state,
spatial, averaging, and scalar errors are respectively below
$6.547\cdot10^{-19}$, $5.455\cdot10^{-19}$, $2.530\cdot10^{-34}$, and
$1.334\cdot10^{-70}$, totaling less than $1.201\cdot10^{-18}$ about an exact
rational center.  The physical time is $2\cdot10^{30}\hbar/\alpha$.
This is a mathematical certificate, not physical feasibility or clock
calibration.  AA1/AA2 are same-author self-review.
Sources: \repo{research/round25/solo/aa1/report.md},
\repo{research/round25/solo/aa2/report.md},
\repo{research/round25/advisor/aa1-gate.json},
\repo{research/round25/advisor/aa2-gate.json}.

\subsubsection{\workstar\ AD1--AD2: null and informative Wilson multipliers}

AD1 tests actual elementary Wilson multiplication.  Forward chooses the yz face
at $O$ and reverse the xy face there; their finite-$q$ collars and constants
remain distinct.  Each $B=\Tr U_f$ has norm two, reference variance one, and
created vector $\psi=B\Omega$ of energy $3\alpha$.  Its square creates an
extra spin-one state: $B\psi=\Omega+\chi_{1,f}$.
For internal cube faces the overlap counts are $r=0,1,4$, so
\eqref{recent:aa-selection} has no zero branch.  The complete exterior charge
argument excludes every external equal-energy image.  Therefore
$\overline A_{q,k}\psi=0=\overline A_{q,k}\Omega$, with participating nonzero
frequencies bounded below by $1/8$.  Applying the ordered scalar identity with
the actual multiplier norm gives, in forward's eight-factor convention,
\begin{align}
 \left|e^{i3\alpha T_q/\hbar}C_q^B(T_q)-1\right|
 &\le24\widehat d_q+4D_k(15z)+48\tau_qM_k(1+3zM_k),\notag\\
 \lim_{q\uparrow1}e^{i3\alpha T_q/\hbar}C_q^B(T_q)&=1.
 \qquad\text{\workstar}
 \label{recent:null-wilson}
\end{align}
The right-column error costs $2b_k$ and the remaining created-vector error
$b_k$; variance one does not make this multiplier norm one.
Reverse's twelve-factor seed uses $(4)_n/n!$ and its own tail
$(4)_{k+1}x^{k+1}/[(k+1)!(1-x)^{k+5}]$, with conservative averaging cost $4b_k$.
The null phase-removed scalar carries no slow-rate information.  Nonzero
variance therefore does not establish an informative probe, and AD1's scale
bridge remains limited.  Its selected forward finite-$q$ example has radius
below $5\cdot10^{-18}$ at the same impractical clock as AA2.
Sources: \repo{research/round26/forward/ad1/report.md},
\repo{research/round26/reverse/ad1/report.md},
\repo{research/round26/skeptic/ad1.md}.

AD2 freezes the actual original six-cycle multipliers
$B_X=\chi_X=\Tr(XZ^{-1})$ and
$B_\Sigma=(\chi_X+\chi_Y)/\sqrt2$.  Their exact multiplication norms are two
and $2\sqrt2$; both have mean zero and variance one.  Their created vectors
are respectively $\phi_X$ and $\psi_+$, but their fourth moments are two and
$5/2$, versus one for the rank observable.  In particular the extra squared
norms of $(B_X^2-I)\Omega$ and $(B_\Sigma^2-I)\Omega$ are one and $3/2$.
Equality of a created vector therefore transfers only the scalar mechanism
proved with the actual multiplier, not the entire operator or its error budget.

The reducing component from AA2 applies to both sources.  Exact polynomial
spectral projectors give for the single X source weights $2/3$ at zero,
$1/8$ at each of $\pm2$, and $1/24$ at each of $\pm2\sqrt3$.
Its moments are $1,0,2,0,16,0,160$.  Consequently, for the full stationary
correlations on the original clock,
\begin{align}
 F_X(z)&=\frac23+\frac14\cos(2\theta)
                   +\frac1{12}\cos(2\sqrt3\theta),\notag\\
 F_\Sigma(z)&=F_\infty(z)\text{ of \eqref{recent:aa-endpoint}},
 \qquad\theta=z/84,\quad0<z\le10^{-6},\notag\\
 \theta^2-\tfrac23\theta^4&\le1-F_X(z)\le\theta^2,\qquad
 |F_\Sigma(z)-1-i\theta|\le2\theta^2.\qquad\text{\workstar}
 \label{recent:wilson-endpoints}
\end{align}
The positive lower quadratic witness follows from
$1-\cos x\ge x^2/2-x^4/24$ and the complete second/fourth moments.  The coherent
sum has a linear imaginary witness.  A constant fitted to X's zero slope fails
its curvature; a single cosine fitted to second moment two predicts fourth
moment four, whereas the actual fourth moment is sixteen.  A single phase fitted
to the coherent mean one predicts second moment one, whereas the actual value
is four.  These are discriminators of named reduced descriptions, not universal
model-identification claims.

For a source $\psi_j=B_j\Omega$ of norm one and $J_j=\norm{B_j}$, the complete
finite-$q$ budget is
\begin{equation}
 \left|e^{iET_q/\hbar}C_q^{B_j}(T_q)
 -\langle\psi_j,e^{is_qQ(q)/96}\psi_j\rangle\right|
 \le6J_j^2\widehat d_q+J_j^2D_k(15z)+(1+J_j)b_k.
 \label{recent:wilson-finite-q}
\end{equation}
The factor $1+J_j$ follows by first replacing the vacuum column and then using
the normalized created source on the remaining column.  Use $J_X^2=4$,
$J_\Sigma^2=8$, and $1+2\sqrt2<4$; rank-norm budgets would undercount state
and spatial errors.  At the AA2 sample parameters, degree-eight exact rational
centers have total radii below $5\cdot10^{-18}$ for X and $10^{-17}$ for the
sum.  Their certified displacements exceed these radii, giving actual finite-$q$
as well as limiting witnesses.  If comparing to the limiting trigonometric
formulas, the coefficient/time correction after \eqref{recent:aa-finite-q}
remains separate.

The Wilson obstruction in V1 is thus resolved by new source-specific observables,
not by falsely approximating the rank operator in the forbidden norm.  The
remaining bridge is empirical or cross-model scale identification with frozen
observables, generators, preparations, and held-out data.  The elementary null
probe, multiple-frequency witnesses, and physical-time budgets establish none
of that matching by themselves.  They also supply no global $q=1$ dynamics,
homogeneous stability theorem, or four-dimensional continuum mass gap.
Sources: \repo{research/round26/forward/ad2/report.md},
\repo{research/round26/reverse/ad2/report.md},
\repo{research/round26/skeptic/ad2.md}.
